\documentclass[../main.tex]{subfiles}
\graphicspath{{\subfix{../../images/}}}

\begin{document}

\section{Installation}

This section will guide you through the installation process of NoTeX and its companion utility.

\subsection{Requirements}

The project requires the following dependencies installed in your system:

\begin{itemize}
    \item CMake (>= 3.20)
    \item C++ compiler (C++20)
    \item \href{https://www.latex-project.org/get/}{LaTeX distribution}, with the following packages:
    \begin{itemize}
        \item \cc{texlive-latex-base}
        \item \cc{texlive-latex-extra}
        \item \cc{texlive-fonts-recommended}
        \item \cc{texlive-fonts-extra}
        \item \cc{texlive-luatex}
        \item \cc{texlive-lang-english}
        \item \cc{texlive-bibtex-extra}
        \item \cc{texlive-science}
        \item \cc{biber}
        \item \cc{latexmk}
    \end{itemize}
    \item \href{https://www.luatex.org/}{LuaLaTeX} compiler: this can be installed with the above \cc{texlive-luatex} package
    \item \href{https://pygments.org/}{Pygments} python package: can be installed with \cc{pip install Pygments} in your preferred environment
\end{itemize}

\subsection{CLI Utility Installation}

\subsubsection{Building from source}

The NoTeX software can be built from source following these steps:

\begin{enumerate}
    \item Clone the repository:
    
        \begin{bash}
            git clone https://github.com/marcotallone/notex.git
        \end{bash}

    \item Go to the project's root:
    
        \begin{bash}
            cd notex
        \end{bash}

    \item Compile using cmake:
    
        \begin{bash}
            mkdir -p build \
            && cmake -S . -B build  -DCMAKE_BUILD_TYPE=Release \
            && cmake --build build --parallel 4
        \end{bash}

        This will produce a unique binary in \cc{./bin/notex}, containing the CLI utility with embedded LaTeX template.

        
    \item \itt{Optionally}, you can install the software system-wise for easier usage:
    
        \begin{bash}
            sudo cmake --install build
        \end{bash}

        This will allow to use \cc{notex} as a normal binary everywhere in your system.

        \begin{warning}
            The system-wise installation might require \cc{sudo} priviledges!
        \end{warning}

    \item Check if your installation was successful by running:

        \begin{bash}
            notex info
            # or, if you skipped step 3:
            # ./bin/notex info
        \end{bash}

        The [`install.sh`](../../install.sh) script provided in the project's repository helps to automate these steps with:

        \begin{bash}
            ./install.sh # for local compilation only
        \end{bash}

        or:

        \begin{bash}
            sudo ./install.sh global # for system-wise installation
        \end{bash}

        \begin{blueinfo}
            Mind that the latter option also installs the LaTeX template system-wise, see below for further details.
        \end{blueinfo}
\end{enumerate}

\subsubsection{Build Options}

When building from source, the CMake compilation offers multiple options that are not set by default:

\begin{itemize}
    \item \cc{DEBUG}: enables debug mode, with extended logging reporting.
    \item \cc{VERBOSE}: enables verbose mode.
    \item \cc{LOGGING}: enables logging reporting. Logs are always stored in the\\ \cc{./log/notex.log} file in the project where \cc{notex} is used.
    \item \cc{TESTS}: enables tests. See [Testing](./Testing.md) for details.
    \item \cc{COVERAGE}: enables coverage report. See [Testing](./Testing.md) for details.
    \item \cc{DOCS}: allows to build project documentation providing the \cc{docs} target. After setting this option with \cc{-DDOCS=ON}, documentation can be built in \cc{./build/docs/} with:
    
    \begin{bash}
        cmake --build build --target docs
    \end{bash}
\end{itemize}

\subsubsection{Uninstall}

You can uninstall a global \cc{notex} installation with the command:

\begin{bash}
    notex uninstall
\end{bash}

Or, equivalently:

\begin{bash}
    notex prune
\end{bash}

A local installation can be analogously uninstalled by providing the path to the project installation:

\begin{bash}
    notex uninstall <path>
\end{bash}

\subsection{LaTeX Template Installation}

After installing the \cc{notex} CLI utility, the LaTeX template can be installed either \bft{globally} (system-wise available) or \bft{locally} (for the single LaTeX project).

\subsubsection{Global Installation}

The template can be globally installed with:

\begin{bash}
    notex install
\end{bash}

Which will install the LaTeX classes and styles, as well as required fonts, into your TeX distribution's \cc{TEXMFHOME} directory. To check which directory this is in your own system, you can run:

\begin{bash}
    kpsewhich -var-value=TEXMFHOME
\end{bash}

\begin{info}
    Pass \cc{--force} to overwrite an existing installation without being asked to confirm.
\end{info}

After a global installation the template can be loaded as:

\begin{latex}
\documentclass[OPTIONS]{notex}
\end{latex}


\subsubsection{Local Installation}

You can install the template locally, for your project only, by running:

\begin{bash}
    notex install <path>
\end{bash}

where \cc{<path>} is the path to your LaTeX project root directory. For instance if the root directory of your project is the current one run:

\begin{bash}
    notex install .
\end{bash}

The local installation will install the template and its classes in the \cc{<path>/settings/} directory, while fonts will be placed in \cc{<path>/fonts/}.

After a local installation the template can be loaded as:

\begin{latex}
\documentclass[OPTIONS]{settings/notex}
\end{latex}

\end{document}
